\captionsetup{type=figure, list=no}
\phantomsection
\begin{lstlisting}[frame=none, basicstyle=\small, breaklines=true, breakatwhitespace=false, breakindent=0pt]
You are a senior software developer with expertise in Java, jOOQ and SQL. Analyze the provided Java class and check for the following database design antipatterns, as defined by Bill Karwin:

- ID Required: Never identify this issue in classes representing views, as views cannot contain primary keys. If the class represents a table, always detect the antipattern, if the name of the primary key is just "id" (case-insensitive). Also detect the issue if a synthetic primary key column exists, even though another unique constraint exists, which is suitable as a primary key (the constraint is on columns, which are virtually immutable by nature), and which does not complicate foreign keys referencing the table too much. Only include the lines of the primary key column definition in the line range, do not include comments or anything else.
- Keyless Entry: A column, which refers to another table, is missing its foreign key. Never identify this issue in classes representing views, as views cannot contain foreign keys. Only report this issue if the "Keys" class provided at the end contains a primary key that this is appropriate for this column to refer to.
- Rounding Errors: Storing fixed-precision values in floating-point type columns, such as FLOAT and REAL, rather than using fixed-precision types like DECIMAL and NUMERIC.
- 31 Flavors: Specifying allowed values in the column definition, i.e. with a CHECK constraint or an ENUM type, rather than using a lookup table. Only include the lines of the column definition in the line range, do not include comments or anything else. Do not report the issue if the CHECK constraint is used to check the value for emptiness or against a range of values (including greater/lesser than comparisons).
- Fear of the Unknown: A special default value, such as an empty string, is used to mark a missing value, rather than NULL, and the special value does not hold a semantic meaning. A column, which can never be NULL in practice (e.g. it has a default value), is marked as NULLABLE.

If the file does not contain any antipatterns, leave the list of occurrences empty.

<analyzed_class>
1: /*
2: * This file is generated by jOOQ.
3: */
4: package edu.java.scrapper.model.jooq.tables;
5: 
6: import edu.java.scrapper.model.jooq.DefaultSchema;
7: import edu.java.scrapper.model.jooq.Keys;
8: import java.util.Arrays;
9: import java.util.List;
10: import javax.annotation.processing.Generated;
11: import org.jetbrains.annotations.NotNull;
12: import org.jetbrains.annotations.Nullable;
13: import org.jooq.Field;
14: import org.jooq.ForeignKey;
15: import org.jooq.Identity;
16: import org.jooq.Name;
17: import org.jooq.Record;
18: import org.jooq.Schema;
19: import org.jooq.Table;
20: import org.jooq.TableField;
21: import org.jooq.TableOptions;
22: import org.jooq.UniqueKey;
23: import org.jooq.impl.DSL;
24: import org.jooq.impl.SQLDataType;
25: import org.jooq.impl.TableImpl;
26: 
27: 
28: /**
29: * This class is generated by jOOQ.
30: */
31: @Generated(
32: value = {
33: "https://www.jooq.org",
34: "jOOQ version:3.18.4"
35: },
36: comments = "This class is generated by jOOQ"
37: )
38: @SuppressWarnings({ "all", "unchecked", "rawtypes" })
39: public class Link extends TableImpl<Record> {
40: 
41: private static final long serialVersionUID = 1L;
42: 
43: /**
44: * The reference instance of <code>LINK</code>
45: */
46: public static final Link LINK = new Link();
47: 
48: /**
49: * The class holding records for this type
50: */
51: @Override
52: @NotNull
53: public Class<Record> getRecordType() {
54: return Record.class;
55: }
56: 
57: /**
58: * The column <code>LINK.ID</code>.
59: */
60: public final TableField<Record, Long> ID = createField(DSL.name("ID"), SQLDataType.BIGINT.nullable(false).identity(true), this, "");
61: 
62: /**
63: * The column <code>LINK.URL</code>.
64: */
65: public final TableField<Record, String> URL = createField(DSL.name("URL"), SQLDataType.VARCHAR(1000000000).nullable(false), this, "");
66: 
67: private Link(Name alias, Table<Record> aliased) {
68: this(alias, aliased, null);
69: }
70: 
71: private Link(Name alias, Table<Record> aliased, Field<?>[] parameters) {
72: super(alias, null, aliased, parameters, DSL.comment(""), TableOptions.table());
73: }
74: 
75: /**
76: * Create an aliased <code>LINK</code> table reference
77: */
78: public Link(String alias) {
79: this(DSL.name(alias), LINK);
80: }
81: 
82: /**
83: * Create an aliased <code>LINK</code> table reference
84: */
85: public Link(Name alias) {
86: this(alias, LINK);
87: }
88: 
89: /**
90: * Create a <code>LINK</code> table reference
91: */
92: public Link() {
93: this(DSL.name("LINK"), null);
94: }
95: 
96: public <O extends Record> Link(Table<O> child, ForeignKey<O, Record> key) {
97: super(child, key, LINK);
98: }
99: 
100: @Override
101: @Nullable
102: public Schema getSchema() {
103: return aliased() ? null : DefaultSchema.DEFAULT_SCHEMA;
104: }
105: 
106: @Override
107: @NotNull
108: public Identity<Record, Long> getIdentity() {
109: return (Identity<Record, Long>) super.getIdentity();
110: }
111: 
112: @Override
113: @NotNull
114: public UniqueKey<Record> getPrimaryKey() {
115: return Keys.CONSTRAINT_2;
116: }
117: 
118: @Override
119: @NotNull
120: public List<UniqueKey<Record>> getUniqueKeys() {
121: return Arrays.asList(Keys.CONSTRAINT_23);
122: }
123: 
124: @Override
125: @NotNull
126: public Link as(String alias) {
127: return new Link(DSL.name(alias), this);
128: }
129: 
130: @Override
131: @NotNull
132: public Link as(Name alias) {
133: return new Link(alias, this);
134: }
135: 
136: @Override
137: @NotNull
138: public Link as(Table<?> alias) {
139: return new Link(alias.getQualifiedName(), this);
140: }
141: 
142: /**
143: * Rename this table
144: */
145: @Override
146: @NotNull
147: public Link rename(String name) {
148: return new Link(DSL.name(name), null);
149: }
150: 
151: /**
152: * Rename this table
153: */
154: @Override
155: @NotNull
156: public Link rename(Name name) {
157: return new Link(name, null);
158: }
159: 
160: /**
161: * Rename this table
162: */
163: @Override
164: @NotNull
165: public Link rename(Table<?> name) {
166: return new Link(name.getQualifiedName(), null);
167: }
168: }
</analyzed_class>

<key_definitions_for_reference>
/*
* This file is generated by jOOQ.
*/
package edu.java.scrapper.model.jooq;

import edu.java.scrapper.model.jooq.tables.GitRepository;
import edu.java.scrapper.model.jooq.tables.Link;
import edu.java.scrapper.model.jooq.tables.StackoverflowQuestion;
import edu.java.scrapper.model.jooq.tables.TgChat;
import edu.java.scrapper.model.jooq.tables.TgChatLink;
import javax.annotation.processing.Generated;
import org.jooq.ForeignKey;
import org.jooq.Record;
import org.jooq.TableField;
import org.jooq.UniqueKey;
import org.jooq.impl.DSL;
import org.jooq.impl.Internal;


/**
* A class modelling foreign key relationships and constraints of tables in the
* default schema.
*/
@Generated(
value = {
"https://www.jooq.org",
"jOOQ version:3.18.4"
},
comments = "This class is generated by jOOQ"
)
@SuppressWarnings({ "all", "unchecked", "rawtypes" })
public class Keys {

// -------------------------------------------------------------------------
// UNIQUE and PRIMARY KEY definitions
// -------------------------------------------------------------------------

public static final UniqueKey<Record> CONSTRAINT_3 = Internal.createUniqueKey(GitRepository.GIT_REPOSITORY, DSL.name("CONSTRAINT_3"), new TableField[] { GitRepository.GIT_REPOSITORY.ID }, true);
public static final UniqueKey<Record> CONSTRAINT_392 = Internal.createUniqueKey(GitRepository.GIT_REPOSITORY, DSL.name("CONSTRAINT_392"), new TableField[] { GitRepository.GIT_REPOSITORY.URN }, true);
public static final UniqueKey<Record> CONSTRAINT_2 = Internal.createUniqueKey(Link.LINK, DSL.name("CONSTRAINT_2"), new TableField[] { Link.LINK.ID }, true);
public static final UniqueKey<Record> CONSTRAINT_23 = Internal.createUniqueKey(Link.LINK, DSL.name("CONSTRAINT_23"), new TableField[] { Link.LINK.URL }, true);
public static final UniqueKey<Record> CONSTRAINT_9 = Internal.createUniqueKey(StackoverflowQuestion.STACKOVERFLOW_QUESTION, DSL.name("CONSTRAINT_9"), new TableField[] { StackoverflowQuestion.STACKOVERFLOW_QUESTION.ID }, true);
public static final UniqueKey<Record> CONSTRAINT_9D3 = Internal.createUniqueKey(StackoverflowQuestion.STACKOVERFLOW_QUESTION, DSL.name("CONSTRAINT_9D3"), new TableField[] { StackoverflowQuestion.STACKOVERFLOW_QUESTION.URN }, true);
public static final UniqueKey<Record> CONSTRAINT_D = Internal.createUniqueKey(TgChat.TG_CHAT, DSL.name("CONSTRAINT_D"), new TableField[] { TgChat.TG_CHAT.ID }, true);
public static final UniqueKey<Record> CONSTRAINT_816 = Internal.createUniqueKey(TgChatLink.TG_CHAT_LINK, DSL.name("CONSTRAINT_816"), new TableField[] { TgChatLink.TG_CHAT_LINK.TG_CHAT_ID, TgChatLink.TG_CHAT_LINK.LINK_ID }, true);

// -------------------------------------------------------------------------
// FOREIGN KEY definitions
// -------------------------------------------------------------------------

public static final ForeignKey<Record, Record> CONSTRAINT_39 = Internal.createForeignKey(GitRepository.GIT_REPOSITORY, DSL.name("CONSTRAINT_39"), new TableField[] { GitRepository.GIT_REPOSITORY.LINK_ID }, Keys.CONSTRAINT_2, new TableField[] { Link.LINK.ID }, true);
public static final ForeignKey<Record, Record> CONSTRAINT_9D = Internal.createForeignKey(StackoverflowQuestion.STACKOVERFLOW_QUESTION, DSL.name("CONSTRAINT_9D"), new TableField[] { StackoverflowQuestion.STACKOVERFLOW_QUESTION.LINK_ID }, Keys.CONSTRAINT_2, new TableField[] { Link.LINK.ID }, true);
public static final ForeignKey<Record, Record> CONSTRAINT_8 = Internal.createForeignKey(TgChatLink.TG_CHAT_LINK, DSL.name("CONSTRAINT_8"), new TableField[] { TgChatLink.TG_CHAT_LINK.TG_CHAT_ID }, Keys.CONSTRAINT_D, new TableField[] { TgChat.TG_CHAT.ID }, true);
public static final ForeignKey<Record, Record> CONSTRAINT_81 = Internal.createForeignKey(TgChatLink.TG_CHAT_LINK, DSL.name("CONSTRAINT_81"), new TableField[] { TgChatLink.TG_CHAT_LINK.LINK_ID }, Keys.CONSTRAINT_2, new TableField[] { Link.LINK.ID }, true);
}
</key_definitions_for_reference>
\end{lstlisting}
\nopagebreak
\begin{minipage}{\linewidth}
    \captionof{figure}{Prompt for detecting design antipatterns using Zero-Shot Prompting.}
\end{minipage}


\captionsetup{type=figure}
\phantomsection
\begin{lstlisting}[frame=none, basicstyle=\small, breaklines=true, breakatwhitespace=false, breakindent=0pt]
You are a senior software developer with expertise in Java, jOOQ and SQL. Analyze the provided Java class and check for the following SQL query antipatterns, as defined by Bill Karwin:

- Poor Man's Search Engine: Usage of LIKE, ILIKE or regular expressions to perform full-text search. Report the issue if it isn't obvious from the method input parameters, whether the patterns contain wildcards used for full-text search. Do not report the issue if LIKE, ILIKE or regex is used for prefix search. Only include the line(s) where the full-text search condition is created in the line range.
- Implicit Columns: A query fetching all columns from a database table. In addition to obvious violations, report cases where jOOQ fetches all columns of a table into records or generated DAOs (located in a package ending with `tables.daos`). Do not report this issue if it occurs within a `fetchCount` or `fetchExists` call. Only include the line(s) where the blind projection is selected in the line range, do not include the rest of the query.
- Fear of the Unknown: Query logic uses a NULLABLE column in a way that produces incorrect results with NULL. Do not report issues that arise from insufficient null-handling in Java code. Also do not report the issue if you're unsure if the column is NULLABLE.

Only identify problems in code, which interacts directly with the jOOQ DSL or generated DAOs (located in a package ending with `tables.daos`). Do not identify problems in code, which interacts with higher level abstractions. In case of multiple consecutive issues, report them separately, even if they are on consecutive lines. If the file does not contain any antipatterns, leave the list of occurrences empty.

<analyzed_class>
1: package edu.java.scrapper.repositories.jooq;
2: 
3: import edu.java.scrapper.model.Link;
4: import edu.java.scrapper.model.TgChatLink;
5: import edu.java.scrapper.repositories.ChatLinkRepository;
6: import java.util.List;
7: import org.jooq.DSLContext;
8: import static edu.java.scrapper.model.jooq.Tables.LINK;
9: import static edu.java.scrapper.model.jooq.Tables.TG_CHAT_LINK;
10: 
11: public class JooqChatLinkRepository implements ChatLinkRepository {
12: 
13: private final DSLContext context;
14: 
15: public JooqChatLinkRepository(DSLContext context) {
16: this.context = context;
17: }
18: 
19: @Override
20: public void addLink(long chatId, long linkId) {
21: context
22: .insertInto(TG_CHAT_LINK, TG_CHAT_LINK.TG_CHAT_ID, TG_CHAT_LINK.LINK_ID)
23: .values(chatId, linkId)
24: .execute();
25: }
26: 
27: @Override
28: public void removeLink(long chatId, long linkId) {
29: context
30: .delete(TG_CHAT_LINK)
31: .where(TG_CHAT_LINK.TG_CHAT_ID.eq(chatId)).and(TG_CHAT_LINK.LINK_ID.eq(linkId))
32: .execute();
33: }
34: 
35: @Override
36: public List<Link> getLinksByChatId(long chatId) {
37: return context
38: .select(LINK.ID, LINK.URL)
39: .from(TG_CHAT_LINK).leftJoin(LINK).on(TG_CHAT_LINK.LINK_ID.eq(LINK.ID))
40: .where(TG_CHAT_LINK.TG_CHAT_ID.eq(chatId)).fetchInto(Link.class);
41: }
42: 
43: @Override
44: public void deleteChatRelatedLinks(long chatId) {
45: context
46: .delete(TG_CHAT_LINK)
47: .where(TG_CHAT_LINK.TG_CHAT_ID.eq(chatId))
48: .execute();
49: }
50: 
51: @Override
52: public boolean existsByChatAndLinkId(long chatId, long linkId) {
53: 
54: List<TgChatLink> tgChatLinks = context
55: .select().from(TG_CHAT_LINK)
56: .where(TG_CHAT_LINK.TG_CHAT_ID.eq(chatId)).and(TG_CHAT_LINK.LINK_ID.eq(linkId))
57: .fetchInto(TgChatLink.class);
58: return !tgChatLinks.isEmpty();
59: }
60: }
</analyzed_class>
\end{lstlisting}
\nopagebreak
\begin{minipage}{\linewidth}
    \captionof{figure}{Prompt for detecting query antipatterns using Zero-Shot Prompting.}
\end{minipage}


\captionsetup{type=figure}
\phantomsection
\begin{lstlisting}[frame=none, basicstyle=\small, breaklines=true, breakatwhitespace=false, breakindent=0pt]
You are a senior software developer with expertise in Java, jOOQ and SQL. Analyze the provided Java class and check for the following database design antipatterns, as defined by Bill Karwin:

- ID Required: Never identify this issue in classes representing views, as views cannot contain primary keys. If the class represents a table, always detect the antipattern, if the name of the primary key is just "id" (case-insensitive). Also detect the issue if a synthetic primary key column exists, even though another unique constraint exists, which is suitable as a primary key (the constraint is on columns, which are virtually immutable by nature), and which does not complicate foreign keys referencing the table too much. Only include the lines of the primary key column definition in the line range, do not include comments or anything else.
- Keyless Entry: A column, which refers to another table, is missing its foreign key. Never identify this issue in classes representing views, as views cannot contain foreign keys. Only report this issue if the "Keys" class provided at the end contains a primary key that this is appropriate for this column to refer to.
- Rounding Errors: Storing fixed-precision values in floating-point type columns, such as FLOAT and REAL, rather than using fixed-precision types like DECIMAL and NUMERIC.
- 31 Flavors: Specifying allowed values in the column definition, i.e. with a CHECK constraint or an ENUM type, rather than using a lookup table. Only include the lines of the column definition in the line range, do not include comments or anything else. Do not report the issue if the CHECK constraint is used to check the value for emptiness or against a range of values (including greater/lesser than comparisons).
- Fear of the Unknown: A special default value, such as an empty string, is used to mark a missing value, rather than NULL, and the special value does not hold a semantic meaning. A column, which can never be NULL in practice (e.g. it has a default value), is marked as NULLABLE.

If the file does not contain any antipatterns, leave the list of occurrences empty.

<analyzed_class>
(same as above)
</analyzed_class>

<key_definitions_for_reference>
(same as above)
</key_definitions_for_reference>

<example>
<input>
...
59: /**
60:  * The column <code>world.city.ID</code>.
61:  */
62: public final TableField<CityRecord, Integer> ID = createField("ID", org.jooq.impl.SQLDataType.INTEGER.nullable(false).identity(true), this, "");
...
</input>
<output>
{
"antipatternName": "ID Required",
"linesRangeStart": 62,
"linesRangeEnd": 62,
"codeFragment": "public final TableField<CityRecord, Integer> ID = createField("ID", org.jooq.impl.SQLDataType.INTEGER.nullable(false).identity(true), this, "");"
"reasoning": "antipattern - the primary key is called ID"
}
</output>
</example>

<example>
<input>
...
59: /**
60:  * The column <code>person.PERSON_ID</code>.
61:  */
62: public final TableField<PersonRecord, Long> PERSON_ID = createField("person_id", org.jooq.impl.SQLDataType.BIGINT.nullable(false).identity(true), this, "");
63: 
64: /**
65:  * The column <code>person.IDENTIFICATION_CODE</code>.
66:  */
67: public final TableField<PersonRecord, String> IDENTIFICATION_CODE = createField("identification_code", org.jooq.impl.SQLDataType.VARCHAR(11).nullable(false), this, "");
...
@Override
public List<UniqueKey<PersonRecord>> getKeys() {
return Arrays.<UniqueKey<PersonRecord>>asList(Keys.KEY_PERSON_IDENTIFICATION_CODE, Keys.KEY_PERSON_PRIMARY);
}
...
Keys.java:
...
public static final UniqueKey<FileRecord> KEY_PERSON_IDENTIFICATION_CODE = Internal.createUniqueKey(Person.PERSON, "KEY_person_identification_code", Person.PERSON.IDENTIFICATION_CODE);
public static final UniqueKey<FileRecord> KEY_PERSON_PRIMARY = Internal.createUniqueKey(Person.PERSON, "KEY_person_PRIMARY", Person.PERSON.PERSON_ID);
...
</input>
<output>
{
"antipatternName": "ID Required",
"linesRangeStart": 62,
"linesRangeEnd": 62,
"codeFragment": "public final TableField<PersonRecord, Long> PERSON_ID = createField("person_id", org.jooq.impl.SQLDataType.BIGINT.nullable(false).identity(true), this, "");"
"reasoning": "antipattern - there is another suitable unique key"
}
</output>
</example>

<example>
<input>
...
59: /**
60:  * The column <code>person.PERSON_ID</code>.
61:  */
62: public final TableField<PersonRecord, Long> PERSON_ID = createField("person_id", org.jooq.impl.SQLDataType.BIGINT.nullable(false).identity(true), this, "");
63: 
64: /**
65:  * The column <code>person.NAME</code>.
66:  */
67: public final TableField<PersonRecord, String> NAME = createField("name", org.jooq.impl.SQLDataType.VARCHAR(32).nullable(false), this, "");
...
@Override
public List<UniqueKey<PersonRecord>> getKeys() {
return Arrays.<UniqueKey<PersonRecord>>asList(Keys.KEY_PERSON_NAME, Keys.KEY_PERSON_PRIMARY);
}
...
Keys.java:
...
public static final UniqueKey<FileRecord> KEY_PERSON_NAME = Internal.createUniqueKey(Person.PERSON, "KEY_person_name", Person.PERSON.NAME);
public static final UniqueKey<FileRecord> KEY_PERSON_PRIMARY = Internal.createUniqueKey(Person.PERSON, "KEY_person_PRIMARY", Person.PERSON.PERSON_ID);
...
</input>
<output>
(nothing - not an antipattern, the alternative key is too volatile)
</output>
</example>

<example>
<input>
...
40: /**
41:  * The column <code>person.PERSON_ID</code>.
42:  */
43: public final TableField<PersonRecord, Long> PERSON_ID = createField("person_id", org.jooq.impl.SQLDataType.BIGINT.nullable(false).identity(true), this, "");
44: 
45: /**
46:  * The column <code>person.FIRST_NAME</code>.
47:  */
48: public final TableField<PersonRecord, String> FIRST_NAME = createField("first_name", org.jooq.impl.SQLDataType.VARCHAR(50).nullable(false), this, "");
49: 
50: /**
51:  * The column <code>person.LAST_NAME</code>.
52:  */
53: public final TableField<PersonRecord, String> LAST_NAME = createField("last_name", org.jooq.impl.SQLDataType.VARCHAR(50).nullable(false), this, "");
54: 
55: /**
56:  * The column <code>person.DATE_OF_BIRTH</code>.
57:  */
58: public final TableField<PersonRecord, LocalDate> DATE_OF_BIRTH = createField("date_of_birth", org.jooq.impl.SQLDataType.LOCALDATE.nullable(false), this, "");
...
@Override
public List<UniqueKey<PersonRecord>> getKeys() {
return Arrays.<UniqueKey<PersonRecord>>asList(Keys.KEY_PERSON_PRIMARY, Keys.KEY_PERSON_NAME_DOB_UNIQUE);
}
...
Keys.java:
...
public static final UniqueKey<PersonRecord> KEY_PERSON_PRIMARY = Internal.createUniqueKey(Person.PERSON, "KEY_person_PRIMARY", Person.PERSON.PERSON_ID);
public static final UniqueKey<PersonRecord> KEY_PERSON_NAME_DOB_UNIQUE = Internal.createUniqueKey(Person.PERSON, "KEY_person_name_dob_unique", Person.PERSON.FIRST_NAME, Person.PERSON.LAST_NAME, Person.PERSON.DATE_OF_BIRTH);
...
</input>
<output>
(nothing - not an antipattern, the alternative key would complicate foreign keys too much)
</output>
</example>

<example>
<input>
...
61: /**
62:  * The column <code>PUBLIC.DEPARTMENT.ORG_ID</code>.
63:  */
64: public final TableField<DepartmentRecord, Long> ORG_ID = createField(DSL.name("ORG_ID"), SQLDataType.BIGINT.nullable(false), this, "");
...
Keys.java:
...
public static final UniqueKey<OrgRecord> CONSTRAINT_1 = Internal.createUniqueKey(Org.ORG, DSL.name("CONSTRAINT_1"), new TableField[] { Org.ORG.ID }, true);
...
(no DEPARTMENT foreign key referencing CONSTRAINT_1 available)
</input>
<output>
{
"antipatternName": "Keyless Entry",
"linesRangeStart": 64,
"linesRangeEnd": 64,
"codeFragment": "public final TableField<DepartmentRecord, Long> ORG_ID = createField(DSL.name("ORG_ID"), SQLDataType.BIGINT.nullable(false), this, "");"
"reasoning": "antipattern - foreign key missing with suitable target available"
}
</output>
</example>

<example>
<input>
...
81: /**
82:  * The column <code>PUBLIC.DEPARTMENT.CREATED_BY</code>.
83:  */
84: public final TableField<DepartmentRecord, Long> CREATED_BY = createField(DSL.name("CREATED_BY"), SQLDataType.BIGINT, this, "");
...
Keys.java:
...
public static final UniqueKey<UserRecord> UX_USER_ID = Internal.createUniqueKey(User.USER, DSL.name("UX_USER_ID"), new TableField[] { User.USER.ID }, true);
...
</input>
<output>
{
"antipatternName": "Keyless Entry",
"linesRangeStart": 84,
"linesRangeEnd": 84,
"codeFragment": "public final TableField<DepartmentRecord, Long> CREATED_BY = createField(DSL.name("CREATED_BY"), SQLDataType.BIGINT, this, "");"
"reasoning": "antipattern - foreign key missing with suitable target available"
}
</output>
</example>

<example>
<input>
...
50: /**
51:  * The column <code>public.scenario.plan_id</code>.
52:  */
53: public final TableField<ScenarioRecord, Long> PLAN_ID = createField(DSL.name("plan_id"), SQLDataType.BIGINT.nullable(false), this, "");
...
Keys.java:
...
(no foreign key on the PLAN_ID column, but also no primary key for a table related to plans available)
...
</input>
<output>
(nothing - not an antipattern, suitable primary key not available)
</output>
</example>

<example>
<input>
...
79: /**
80:  * The column <code>world.country.SurfaceArea</code>.
81:  */
82: public final TableField<CountryRecord, Double> SURFACEAREA = createField("SurfaceArea", org.jooq.impl.SQLDataType.FLOAT.nullable(false).defaultValue(org.jooq.impl.DSL.inline("0.00", org.jooq.impl.SQLDataType.FLOAT)), this, "");
...
</input>
<output>
{
"antipatternName": "Rounding Errors",
"linesRangeStart": 82,
"linesRangeEnd": 82,
"codeFragment": "public final TableField<CountryRecord, Double> SURFACEAREA = createField("SurfaceArea", org.jooq.impl.SQLDataType.FLOAT.nullable(false).defaultValue(org.jooq.impl.DSL.inline("0.00", org.jooq.impl.SQLDataType.FLOAT)), this, "");"
"reasoning": "antipattern - pretty much every use of an imprecise floating point type is incorrect"
}
</output>
</example>

<example>
<input>
...
121: /**
122:  * The column <code>public.film.rating</code>.
123:  */
124: public final TableField<FilmRecord, MpaaRating> RATING = createField(DSL.name("rating"), SQLDataType.VARCHAR.defaultValue(DSL.field(DSL.raw("'G'::mpaa_rating"), SQLDataType.VARCHAR)).asEnumDataType(MpaaRating.class), this, "");
...
</input>
<output>
{
"antipatternName": "31 Flavors",
"linesRangeStart": 124,
"linesRangeEnd": 124,
"codeFragment": "public final TableField<FilmRecord, MpaaRating> RATING = createField(DSL.name("rating"), SQLDataType.VARCHAR.defaultValue(DSL.field(DSL.raw("'G'::mpaa_rating"), SQLDataType.VARCHAR)).asEnumDataType(MpaaRating.class), this, "");"
"reasoning": "antipattern - use of an ENUM to restrict value options"
}
</output>
</example>

<example>
<input>
...
68: /**
69:  * The column <code>pasto.tipo</code>.
70:  */
71: public final TableField<PastoRecord, String> TIPO = createField(DSL.name("tipo"), SQLDataType.CLOB.nullable(false), this, "");
...
126: @Override
127: public List<Check<PastoRecord>> getChecks() {
128: return Arrays.asList(
129: Internal.createCheck(this, DSL.name(""), "tipo in ('colazione', 'pranzo', 'cena', 'spuntino')", true)
130: );
131: }
...
</input>
<output>
{
"antipatternName": "31 Flavors",
"linesRangeStart": 124,
"linesRangeEnd": 124,
"codeFragment": "public final TableField<PastoRecord, String> TIPO = createField(DSL.name("tipo"), SQLDataType.CLOB.nullable(false), this, "");"
"reasoning": "antipattern - use of an CHECK constraint to restrict value options"
}
</output>
</example>

<example>
<input>
...
72: /**
73:  * The column <code>public.dragon.name</code>.
74:  */
75: public final TableField<DragonRecord, String> NAME = createField(DSL.name("name"), SQLDataType.VARCHAR.nullable(false), this, "");
...
84: /**
85:  * The column <code>public.dragon.age</code>.
86:  */
87: public final TableField<DragonRecord, Integer> AGE = createField(DSL.name("age"), SQLDataType.INTEGER.nullable(false), this, "");
...
215: @Override
216: public List<Check<DragonRecord>> getChecks() {
217: return Arrays.asList(
218: Internal.createCheck(this, DSL.name("dragon_age_check"), "((age > 0))", true),
219: Internal.createCheck(this, DSL.name("dragon_name_check"), "(((name)::text <> ''::text))", true)
220: );
221: }
...
</input>
<output>
(nothing - neither of the CHECK constraints validate the value against a predetermined set of values (emptiness check doesn't count either))
</output>
</example>

<example>
<input>
...
59: /**
60:  * The column <code>world.country.Code</code>.
61:  */
62: public final TableField<CountryRecord, String> CODE = createField("Code", org.jooq.impl.SQLDataType.CHAR(3).nullable(false).defaultValue(org.jooq.impl.DSL.inline("", org.jooq.impl.SQLDataType.CHAR)), this, "");
...
</input>
<output>
{
"antipatternName": "Fear of the Unknown",
"linesRangeStart": 62,
"linesRangeEnd": 62,
"codeFragment": "public final TableField<CountryRecord, String> CODE = createField("Code", org.jooq.impl.SQLDataType.CHAR(3).nullable(false).defaultValue(org.jooq.impl.DSL.inline("", org.jooq.impl.SQLDataType.CHAR)), this, "");"
"reasoning": "antipattern - use of a special non-business value to represent a missing value, rather than null"
}
</output>
</example>

<example>
<input>
...
77: /**
78:  * The column <code>public.link.updatetime</code>.
79:  */
80: public final TableField<LinkRecord, LocalDateTime> UPDATETIME = createField(DSL.name("updatetime"), SQLDataType.LOCALDATETIME(6).defaultValue(DSL.field("now()", SQLDataType.LOCALDATETIME)), this, "");
...
</input>
<output>
{
"antipatternName": "Fear of the Unknown",
"linesRangeStart": 80,
"linesRangeEnd": 80,
"codeFragment": "public final TableField<LinkRecord, LocalDateTime> UPDATETIME = createField(DSL.name("updatetime"), SQLDataType.LOCALDATETIME(6).defaultValue(DSL.field("now()", SQLDataType.LOCALDATETIME)), this, "");"
"reasoning": "antipattern - a column, which can not be null in practice, is nullable"
}
</output>
</example>

<example>
<input>
...
93: /**
94:  * The column <code>texera_db.workflow_executions.name</code>.
95:  */
96: public final TableField<WorkflowExecutionsRecord, String> NAME = createField(DSL.name("name"), org.jooq.impl.SQLDataType.VARCHAR(128).nullable(false).defaultValue(org.jooq.impl.DSL.inline("Untitled Execution", org.jooq.impl.SQLDataType.VARCHAR)), this, "");
...
</input>
<output>
(nothing - the default value has a semantic meaning)
</output>
</example>
\end{lstlisting}
\nopagebreak
\begin{minipage}{\linewidth}
    \captionof{figure}{Prompt for detecting design antipatterns using Few-Shot Prompting.}
\end{minipage}


\captionsetup{type=figure}
\phantomsection
\begin{lstlisting}[frame=none, basicstyle=\small, breaklines=true, breakatwhitespace=false, breakindent=0pt]
You are a senior software developer with expertise in Java, jOOQ and SQL. Analyze the provided Java class and check for the following SQL query antipatterns, as defined by Bill Karwin:

- Poor Man's Search Engine: Usage of LIKE, ILIKE or regular expressions to perform full-text search. Report the issue if it isn't obvious from the method input parameters, whether the patterns contain wildcards used for full-text search. Do not report the issue if LIKE, ILIKE or regex is used for prefix search. Only include the line(s) where the full-text search condition is created in the line range.
- Implicit Columns: A query fetching all columns from a database table. In addition to obvious violations, report cases where jOOQ fetches all columns of a table into records or generated DAOs (located in a package ending with `tables.daos`). Do not report this issue if it occurs within a `fetchCount` or `fetchExists` call. Only include the line(s) where the blind projection is selected in the line range, do not include the rest of the query.
- Fear of the Unknown: Query logic uses a NULLABLE column in a way that produces incorrect results with NULL. Do not report issues that arise from insufficient null-handling in Java code. Also do not report the issue if you're unsure if the column is NULLABLE.

Only identify problems in code, which interacts directly with the jOOQ DSL or generated DAOs (located in a package ending with `tables.daos`). Do not identify problems in code, which interacts with higher level abstractions. In case of multiple consecutive issues, report them separately, even if they are on consecutive lines. If the file does not contain any antipatterns, leave the list of occurrences empty.

<analyzed_class>
(same as above)
</analyzed_class>

<example>
<input>
...
20: @Transactional(readOnly = true)
21: public List<String> getTagSuggestionsForImage(String category, String tag, UUID imageId) {
22: UUID repoId = getRepoId(imageId);
23:
24: return
25: dsl.select(Tables.IMAGE_TAGS.TAG)
26: .from(Tables.IMAGE_TAGS)
27: .where(Tables.IMAGE_TAGS.REPOSITORY_ID.eq(repoId))
28: .and(Tables.IMAGE_TAGS.TAG.likeIgnoreCase("%"+tag+"%"))
29: .and(Tables.IMAGE_TAGS.TAG_CATEGORY.likeIgnoreCase("%"+category+"%"))
30: .and(Tables.IMAGE_TAGS.IMAGE_ID.notEqual(imageId))
31: .limit(5)
32: .fetch(Tables.IMAGE_TAGS.TAG);
33: }
...
</input>
<output>
{
"antipatternName": "Poor Man's Search Engine",
"linesRangeStart": 28,
"linesRangeEnd": 28,
"codeFragment": ".and(Tables.IMAGE_TAGS.TAG.likeIgnoreCase("%"+tag+"%"))",
"reasoning": "antipattern - use of ILIKE for full-text search"
}, {
"antipatternName": "Poor Man's Search Engine",
"linesRangeStart": 29,
"linesRangeEnd": 29,
"codeFragment": ".and(Tables.IMAGE_TAGS.TAG_CATEGORY.likeIgnoreCase("%"+category+"%"))",
"reasoning": "antipattern - use of ILIKE for full-text search"
}
</output>
</example>

<example>
<input>
(tag and category are passed in from outside the class)
...
20: @Transactional(readOnly = true)
21: public List<String> getTagSuggestionsForImage(String category, String tag, UUID imageId) {
22: UUID repoId = getRepoId(imageId);
23:
24: return
25: dsl.select(Tables.IMAGE_TAGS.TAG)
26: .from(Tables.IMAGE_TAGS)
27: .where(Tables.IMAGE_TAGS.REPOSITORY_ID.eq(repoId))
28: .and(Tables.IMAGE_TAGS.TAG.likeIgnoreCase(tag))
29: .and(Tables.IMAGE_TAGS.TAG_CATEGORY.likeIgnoreCase(category))
30: .and(Tables.IMAGE_TAGS.IMAGE_ID.notEqual(imageId))
31: .limit(5)
32: .fetch(Tables.IMAGE_TAGS.TAG);
33: }
...
</input>
<output>
{
"antipatternName": "Poor Man's Search Engine",
"linesRangeStart": 28,
"linesRangeEnd": 28,
"codeFragment": ".and(Tables.IMAGE_TAGS.TAG.likeIgnoreCase("%"+tag+"%"))",
"reasoning": "antipattern - use of ILIKE, not able to figure out from the class if it's used for full-text search"
}, {
"antipatternName": "Poor Man's Search Engine",
"linesRangeStart": 29,
"linesRangeEnd": 29,
"codeFragment": ".and(Tables.IMAGE_TAGS.TAG_CATEGORY.likeIgnoreCase("%"+category+"%"))",
"reasoning": "antipattern - use of ILIKE, not able to figure out from the class if it's used for full-text search"
}
</output>
</example>

<example>
<input>
62: @Override
63: public List<Link> findByIdIn(List<Long> ids) {
64: return context
65: .select().from(LINK).where(LINK.ID.in(ids)).fetchInto(Link.class);
66: }
</input>
<output>
{
"antipatternName": "Implicit Columns",
"linesRangeStart": 65,
"linesRangeEnd": 65,
"codeFragment": ".select().from(LINK).where(LINK.ID.in(ids)).fetchInto(Link.class);",
"reasoning": "antipattern - select().from(LINK) fetches all columns from the table"
}
</output>
</example>

<example>
<input>
65: private Void preValidateParams(ReplicationParam.Duplicate duplicate) {
66: boolean exists = dslContext.fetchExists(dslContext.selectFrom(IMAGE).where(IMAGE.ID.eq(duplicate.imageId())));
67: 
68: if (!exists) {
69: throw new ImageReplicationUserException("Image with ID " + duplicate.imageId() + " does not exist");
70: }
71: 
72: return null;
73: }
</input>
<output>
(nothing - the selectFrom, which by itself would fetch a record with all columns, is wrapped in a fetchExists)
</output>
</example>

<example>
<input>
...
import com.example.onboardingservice.jooq.tables.daos.OnboardingDocumentsDao;
...
private final OnboardingDocumentsDao documentsDao;
...
47: @Override
48: public List<DocumentDto> getDocuments() {
49: return documentsDao.findAll().stream()
50: .map(it -> DocumentDto.of(it.getKey(), it.getFilename()))
51: .toList();
52: }
...
</input>
<output>
{
"antipatternName": "Implicit Columns",
"linesRangeStart": 49,
"linesRangeEnd": 49,
"codeFragment": "documentsDao.findAll()",
"reasoning": "antipattern - generated DAOs fetch all columns from the table"
}
</output>
</example>

<example>
<input>
...
import com.example.onboardingservice.business.OnboardingDocumentsDao;
...
private final OnboardingDocumentsDao documentsDao;
...
47: @Override
48: public List<DocumentDto> getDocuments() {
49: return documentsDao.findAll().stream()
50: .map(it -> DocumentDto.of(it.getKey(), it.getFilename()))
51: .toList();
52: }
...
</input>
<output>
(nothing - the DAO is not generated and is analysed separately)
</output>
</example>

<example>
<input>
...
34: public Author afficherAuteur(Integer id){
35: return dslContext.selectFrom(AUTHOR).where(AUTHOR.ID.eq(id)).fetchSingleInto(Author.class);
36: }
...
50: public void modifierAuteur(AuthorVo author){
51: Author _author = this.afficherAuteur(author.getId().intValue());
52:
53: dslContext.update(AUTHOR)
54: .set(AUTHOR.FIRST_NAME,author.getFirstName())
55: .set(AUTHOR.LAST_NAME, author.getLastName())
56: .where(AUTHOR.ID.eq(author.getId().intValue())).execute();
57: }
...
</input>
<output>
(nothing - while there is a potential nullability issue here, it's caused by the Java code)
</output>
</example>

<example>
<input>
...
72: private Condition getIdCondition(UUID id) {
73: if (id == null) {
74: return TASK.TASK_TYPE_ID.eq(TASK.TASK_TYPE_ID); // 1==1
75: } else {
76: return TASK.TASK_TYPE_ID.eq(id);
77: }
78: }
...
</input>
<output>
(nothing - assume that TASK_TYPE_ID is non-nullable)
</output>
</example>
\end{lstlisting}
\nopagebreak
\begin{minipage}{\linewidth}
    \captionof{figure}{Prompt for detecting query antipatterns using Few-Shot Prompting.}
\end{minipage}


\captionsetup{type=figure}
\phantomsection
\begin{lstlisting}[frame=none, basicstyle=\small, breaklines=true, breakatwhitespace=false, breakindent=0pt]
You are a senior software developer with expertise in Java, jOOQ and SQL. Analyze the provided Java class and check for the following database design antipatterns, as defined by Bill Karwin:

- ID Required: Never identify this issue in classes representing views, as views cannot contain primary keys. If the class represents a table, always detect the antipattern, if the name of the primary key is just "id" (case-insensitive). Also detect the issue if a synthetic primary key column exists, even though another unique constraint exists, which is suitable as a primary key (the constraint is on columns, which are virtually immutable by nature), and which does not complicate foreign keys referencing the table too much. Only include the lines of the primary key column definition in the line range, do not include comments or anything else.
- Keyless Entry: A column, which refers to another table, is missing its foreign key. Never identify this issue in classes representing views, as views cannot contain foreign keys. Only report this issue if the "Keys" class provided at the end contains a primary key that this is appropriate for this column to refer to.
- Rounding Errors: Storing fixed-precision values in floating-point type columns, such as FLOAT and REAL, rather than using fixed-precision types like DECIMAL and NUMERIC.
- 31 Flavors: Specifying allowed values in the column definition, i.e. with a CHECK constraint or an ENUM type, rather than using a lookup table. Only include the lines of the column definition in the line range, do not include comments or anything else. Do not report the issue if the CHECK constraint is used to check the value for emptiness or against a range of values (including greater/lesser than comparisons).
- Fear of the Unknown: A special default value, such as an empty string, is used to mark a missing value, rather than NULL, and the special value does not hold a semantic meaning. A column, which can never be NULL in practice (e.g. it has a default value), is marked as NULLABLE.

If the file does not contain any antipatterns, leave the list of occurrences empty.

<analyzed_class>
(same as above)
</analyzed_class>

<key_definitions_for_reference>
(same as above)
</key_definitions_for_reference>

Let's think step by step.
\end{lstlisting}
\nopagebreak
\begin{minipage}{\linewidth}
    \captionof{figure}{Prompt for detecting design antipatterns using Chain-of-Thought Prompting.}
\end{minipage}


\captionsetup{type=figure}
\phantomsection
\begin{lstlisting}[frame=none, basicstyle=\small, breaklines=true, breakatwhitespace=false, breakindent=0pt]
You are a senior software developer with expertise in Java, jOOQ and SQL. Analyze the provided Java class and check for the following SQL query antipatterns, as defined by Bill Karwin:

- Poor Man's Search Engine: Usage of LIKE, ILIKE or regular expressions to perform full-text search. Report the issue if it isn't obvious from the method input parameters, whether the patterns contain wildcards used for full-text search. Do not report the issue if LIKE, ILIKE or regex is used for prefix search. Only include the line(s) where the full-text search condition is created in the line range.
- Implicit Columns: A query fetching all columns from a database table. In addition to obvious violations, report cases where jOOQ fetches all columns of a table into records or generated DAOs (located in a package ending with `tables.daos`). Do not report this issue if it occurs within a `fetchCount` or `fetchExists` call. Only include the line(s) where the blind projection is selected in the line range, do not include the rest of the query.
- Fear of the Unknown: Query logic uses a NULLABLE column in a way that produces incorrect results with NULL. Do not report issues that arise from insufficient null-handling in Java code. Also do not report the issue if you're unsure if the column is NULLABLE.

Only identify problems in code, which interacts directly with the jOOQ DSL or generated DAOs (located in a package ending with `tables.daos`). Do not identify problems in code, which interacts with higher level abstractions. In case of multiple consecutive issues, report them separately, even if they are on consecutive lines. If the file does not contain any antipatterns, leave the list of occurrences empty.

<analyzed_class>
(same as above)
</analyzed_class>

Let's think step by step.
\end{lstlisting}
\nopagebreak
\begin{minipage}{\linewidth}
    \captionof{figure}{Prompt for detecting query antipatterns using Chain-of-Thought Prompting.}
\end{minipage}


\captionsetup{type=figure}
\phantomsection
\begin{lstlisting}[frame=none, basicstyle=\small, breaklines=true, breakatwhitespace=false, breakindent=0pt]
Simulate exactly 3 experts working in parallel on the same complete task.

Rules:
1. Every expert must inspect the entire Java class and all target antipatterns.
2. Experts must not split the work by antipattern, method, or line range.
3. On each round, all remaining experts perform the same kind of step on the same full task.
4. If an expert changes their mind, they may drop out, but only after first attempting the full task independently.
5. After the rounds, produce a final consensus based only on issues agreed by the remaining experts.

The task is: Analyze the provided Java class and check for the following database design antipatterns, as defined by Bill Karwin:

- ID Required: Never identify this issue in classes representing views, as views cannot contain primary keys. If the class represents a table, always detect the antipattern, if the name of the primary key is just "id" (case-insensitive). Also detect the issue if a synthetic primary key column exists, even though another unique constraint exists, which is suitable as a primary key (the constraint is on columns, which are virtually immutable by nature), and which does not complicate foreign keys referencing the table too much. Only include the lines of the primary key column definition in the line range, do not include comments or anything else.
- Keyless Entry: A column, which refers to another table, is missing its foreign key. Never identify this issue in classes representing views, as views cannot contain foreign keys. Only report this issue if the "Keys" class provided at the end contains a primary key that this is appropriate for this column to refer to.
- Rounding Errors: Storing fixed-precision values in floating-point type columns, such as FLOAT and REAL, rather than using fixed-precision types like DECIMAL and NUMERIC.
- 31 Flavors: Specifying allowed values in the column definition, i.e. with a CHECK constraint or an ENUM type, rather than using a lookup table. Only include the lines of the column definition in the line range, do not include comments or anything else. Do not report the issue if the CHECK constraint is used to check the value for emptiness or against a range of values (including greater/lesser than comparisons).
- Fear of the Unknown: A special default value, such as an empty string, is used to mark a missing value, rather than NULL, and the special value does not hold a semantic meaning. A column, which can never be NULL in practice (e.g. it has a default value), is marked as NULLABLE.

If the file does not contain any antipatterns, leave the list of occurrences empty.

<analyzed_class>
(same as above)
</analyzed_class>

<key_definitions_for_reference>
(same as above)
</key_definitions_for_reference>
\end{lstlisting}
\nopagebreak
\begin{minipage}{\linewidth}
    \captionof{figure}{Prompt for detecting design antipatterns using Tree-of-Thought Prompting.}
\end{minipage}


\captionsetup{type=figure}
\phantomsection
\begin{lstlisting}[frame=none, basicstyle=\small, breaklines=true, breakatwhitespace=false, breakindent=0pt]
Simulate exactly 3 experts working in parallel on the same complete task.

Rules:
1. Every expert must inspect the entire Java class and all target antipatterns.
2. Experts must not split the work by antipattern, method, or line range.
3. On each round, all remaining experts perform the same kind of step on the same full task.
4. If an expert changes their mind, they may drop out, but only after first attempting the full task independently.
5. After the rounds, produce a final consensus based only on issues agreed by the remaining experts.

The task is: Analyze the provided Java class and check for the following SQL query antipatterns, as defined by Bill Karwin:

- Poor Man's Search Engine: Usage of LIKE, ILIKE or regular expressions to perform full-text search. Report the issue if it isn't obvious from the method input parameters, whether the patterns contain wildcards used for full-text search. Do not report the issue if LIKE, ILIKE or regex is used for prefix search. Only include the line(s) where the full-text search condition is created in the line range.
- Implicit Columns: A query fetching all columns from a database table. In addition to obvious violations, report cases where jOOQ fetches all columns of a table into records or generated DAOs (located in a package ending with `tables.daos`). Do not report this issue if it occurs within a `fetchCount` or `fetchExists` call. Only include the line(s) where the blind projection is selected in the line range, do not include the rest of the query.
- Fear of the Unknown: Query logic uses a NULLABLE column in a way that produces incorrect results with NULL. Do not report issues that arise from insufficient null-handling in Java code. Also do not report the issue if you're unsure if the column is NULLABLE.

Only identify problems in code, which interacts directly with the jOOQ DSL or generated DAOs (located in a package ending with `tables.daos`). Do not identify problems in code, which interacts with higher level abstractions. In case of multiple consecutive issues, report them separately, even if they are on consecutive lines. If the file does not contain any antipatterns, leave the list of occurrences empty.

<analyzed_class>
(same as above)
</analyzed_class>
\end{lstlisting}
\nopagebreak
\begin{minipage}{\linewidth}
    \captionof{figure}{Prompt for detecting query antipatterns using Tree-of-Thought Prompting.}
\end{minipage}
